%_________________
\section{Distribusi normal}
\label{normalDist}

\index{distribution!normal|(}
\index{normal distribution|(}

Di antara semua distribusi yang kita jumpai dalam praktik,
ada satu yang jauh lebih umum daripada yang lain.
Kurva lonceng yang simetris dan unimodal muncul di mana-mana
dalam statistika.
Bahkan, kurva ini begitu umum sehingga orang sering menyebutnya
\termsub{kurva normal}{normal distribution} atau
\termsub{distribusi normal}{normal distribution}\index{distribution!normal|textbf}%
,\footnote{Distribusi ini juga diperkenalkan sebagai distribusi Gaussian,
  dari nama Carl Friedrich Gauss, orang pertama yang memformalkan
  bentuk matematisnya.}
seperti yang ditampilkan pada Gambar~\ref{simpleNormal}.
Variabel seperti skor SAT dan tinggi badan laki-laki dewasa di AS
mendekati distribusi normal.

\begin{figure}[h]
  \centering
  \Figure[Ditampilkan sebuah kurva berbentuk lonceng yang simetris terhadap pusatnya. Inilah distribusi normal. Dari kiri, kurva bermula rendah dan perlahan terangkat dari sumbu horizontal, lalu menanjak lebih curam sebelum kenaikannya melambat dan kurva mendatar di puncak. Dari puncak, kurva mula-mula menurun perlahan, kemudian lebih curam, lalu berangsur mendatar saat mendekati sumbu horizontal. Bentuk distribusi normal yang menyerupai lonceng ini dimiliki banyak distribusi yang akan kita jumpai di sepanjang buku. Selanjutnya, ingatlah bentuk lonceng ini setiap kali distribusi normal dibahas.]{0.5}{simpleNormal}
  \caption{Sebuah kurva normal.}
  \label{simpleNormal}
\end{figure}

\begin{onebox}{Fakta tentang distribusi normal}
  Banyak variabel hampir mengikuti distribusi normal, tetapi tidak ada
  yang tepat mengikuti distribusi normal.
  Jadi, meskipun tidak sempurna untuk satu pun masalah tertentu,
  distribusi normal sangat berguna untuk beragam masalah.
  Kita akan menggunakannya dalam eksplorasi data dan untuk menyelesaikan
  berbagai masalah penting dalam statistika.
\end{onebox}


\subsection{Model distribusi normal}

\termsub{Distribusi normal}{normal distribution} selalu menggambarkan
kurva berbentuk lonceng yang simetris dan unimodal.
Namun, tampilan kurva-kurva tersebut dapat berbeda bergantung pada
rincian modelnya.
Secara khusus, model distribusi normal dapat disesuaikan dengan
dua parameter: rata-rata dan simpangan baku.
Seperti yang mungkin Anda duga, perubahan rata-rata menggeser
kurva lonceng ke kiri atau ke kanan, sedangkan perubahan simpangan baku
meregangkan atau menyempitkan kurva.
Panel kiri Gambar~\ref{twoSampleNormals} menampilkan distribusi normal
dengan rata-rata $0$ dan simpangan baku $1$, sedangkan panel kanan
menampilkan distribusi normal dengan rata-rata $19$ dan simpangan baku $4$.
Gambar~\ref{twoSampleNormalsStacked} menampilkan kedua distribusi ini
pada sumbu yang sama.

\begin{figure}[h]
  \centering
  \Figure[Ditampilkan dua distribusi normal. Distribusi pertama berpusat di 0 dan memiliki simpangan baku 1; kedua ekornya praktis tidak dapat dibedakan dari tinggi 0 untuk nilai yang lebih kecil dari -3 atau lebih besar dari positif 3. Distribusi normal kedua berpusat di 19 dan memiliki simpangan baku 4; tinggi distribusinya tidak dapat dibedakan dari 0 ketika jaraknya dari rata-rata lebih dari 3 simpangan baku.]{0.7}{twoSampleNormals}
  \caption{Kedua kurva merepresentasikan distribusi normal.
      Namun, pusat dan sebarannya berbeda.}
  \label{twoSampleNormals}
\end{figure}

\begin{figure}[h]
  \centering
  \Figure[Ditampilkan dua distribusi normal pada grafik yang sama. Distribusi pertama memiliki rata-rata 0 dan simpangan baku 1. Distribusi kedua memiliki rata-rata 19 dan simpangan baku 4. Salah satu sifat penting yang terlihat pada grafik ialah bahwa, karena setiap distribusi harus memiliki luas 1, distribusi normal dengan simpangan baku 1 tampak jauh lebih sempit sekaligus jauh lebih tinggi daripada distribusi kedua yang memiliki simpangan baku 4.]{0.6}{twoSampleNormalsStacked}
  \caption{Distribusi normal pada
      Gambar~\ref{twoSampleNormals}, tetapi digambar bersama-sama
      dengan skala yang sama.}
  \label{twoSampleNormalsStacked}
\end{figure}

Jika suatu distribusi normal memiliki rata-rata $\mu$ dan simpangan baku
$\sigma$, kita dapat menuliskan distribusinya sebagai $N(\mu, \sigma)$.
Kedua distribusi pada Gambar~\ref{twoSampleNormalsStacked}
dapat dituliskan sebagai
\begin{align*}
N(\mu=0,\sigma=1)
  \quad \text{dan} \quad
  N(\mu=19,\sigma=4)
\end{align*}
Karena rata-rata dan simpangan baku menggambarkan suatu distribusi normal
secara tepat, keduanya disebut
\termsub{parameter}{parameter} distribusi tersebut.
Distribusi normal dengan rata-rata $\mu = 0$ dan
simpangan baku $\sigma = 1$ disebut
\termsub{distribusi normal baku}{standard normal distribution}%
\index{normal distribution!standard|textbf}.

\begin{exercisewrap}
\begin{nexercise}
Tuliskan notasi ringkas untuk distribusi normal
dengan\footnotemark{} \\
%\begin{enumerate}[(a)]
%\setlength{\itemsep}{0mm}
%\item
(a)
    rata-rata~5 dan simpangan baku~3, \\
%\item
(b)
    rata-rata~-100 dan simpangan baku~10, serta \\
%\item
(c)
    rata-rata~2 dan simpangan baku~9.
%\end{enumerate}
\end{nexercise}
\end{exercisewrap}
\footnotetext{(a)~$N(\mu=5,\sigma=3)$.
  (b)~$N(\mu=-100, \sigma=10)$.
  (c)~$N(\mu=2, \sigma=9)$.}


\subsection{Standardisasi dengan skor-Z}

\noindent%
Kita sering ingin menempatkan data pada skala baku agar
perbandingan menjadi lebih masuk akal.

\newcommand{\satmean}{1100}
\newcommand{\satsd}{200}
\newcommand{\actmean}{21}
\newcommand{\actsd}{6}
\newcommand{\annsatscore}{1300}
\newcommand{\annsatzscore}{1}
\newcommand{\tomsatscore}{24}
\newcommand{\tomsatzscore}{0.5}

\begin{examplewrap}
\begin{nexample}{Tabel~\vref{satACTstats} menampilkan rata-rata
    dan simpangan baku skor total pada SAT dan ACT.
    Distribusi skor SAT dan ACT sama-sama hampir normal.
    Misalkan Ann memperoleh skor \annsatscore{} pada SAT dan Tom
    memperoleh skor \tomsatscore{} pada ACT.
    Siapa yang memperoleh hasil relatif lebih baik?}
  \label{actSAT}%
  Kita menggunakan simpangan baku sebagai panduan.
  Skor SAT Ann berada \annsatzscore{} simpangan baku di atas rata-rata:
  $\satmean{} + \satsd{} = \annsatscore{}$.
  Skor ACT Tom berada \tomsatzscore{} simpangan baku di atas rata-rata:
  $\actmean{} + \tomsatzscore{} \times \actsd{} = \tomsatscore{}$.
  Pada Gambar~\ref{satActNormals}, kita dapat melihat bahwa,
  dibandingkan dengan peserta lain, Ann cenderung berprestasi lebih baik
  daripada Tom; jadi skor relatif Ann lebih baik.
\end{nexample}
\end{examplewrap}

\begin{figure}[h]
\centering
\begin{tabular}{l r r}
  \hline
  & SAT & ACT \\
  \hline
  Rata-rata \hspace{0.3cm} & \satmean{} & \actmean{} \\
  Simpangan baku & \satsd{} & \actsd{} \\
  \hline
\end{tabular}
\caption{Rata-rata dan simpangan baku untuk SAT dan ACT.}
\label{satACTstats}
\end{figure}

\begin{figure}
  \centering
  \Figure[Skor Ann dan Tom ditampilkan pada distribusi SAT dan ACT, yang masing-masing digambarkan sebagai distribusi normal. Distribusi SAT memiliki rata-rata 1100 dan simpangan baku 200, sedangkan distribusi ACT memiliki rata-rata 21 dan simpangan baku 6. Skor SAT Ann adalah 1300 dan skor ACT Tom adalah 24. Berdasarkan posisi mereka pada grafik masing-masing, tampak bahwa nilai relatif Ann pada distribusi SAT lebih tinggi daripada nilai relatif Tom pada distribusi ACT.]{0.6}{satActNormals}
  \caption{Skor Ann dan Tom pada distribusi SAT dan ACT.}
  \label{satActNormals}
\end{figure}

\begingroup
\looseness=-1
Contoh~\ref{actSAT} menggunakan teknik standardisasi yang disebut
skor-Z. Metode ini paling sering digunakan untuk pengamatan yang
hampir normal, tetapi dapat digunakan pada distribusi apa pun.
\termsub{Skor-Z}{Z-score}\index{Z@$Z$} suatu pengamatan menyatakan
berapa simpangan baku letak pengamatan tersebut
di atas atau di bawah rata-rata.
Jika pengamatan berada satu simpangan baku di atas rata-rata,
skor-Z-nya adalah~1.
Jika pengamatan berada 1.5 simpangan baku \emph{di bawah} rata-rata,
skor-Z-nya adalah -1.5.
Jika $x$ merupakan pengamatan dari distribusi $N(\mu, \sigma)$,
secara matematis kita mendefinisikan skor-Z sebagai
\par
\endgroup
\begin{align*}
Z = \frac{x - \mu}{\sigma}
\end{align*}
Dengan $\mu_{SAT} = \satmean{}$, $\sigma_{SAT} = \satsd{}$,
dan $x_{_{\text{Ann}}} = \annsatscore{}$, kita memperoleh skor-Z Ann:
\begin{align*}
Z_{_{\text{Ann}}}
  = \frac{x_{_{\text{Ann}}} - \mu_{_{\text{SAT}}}}
      {\sigma_{_{\text{SAT}}}}
  = \frac{\annsatscore{} - \satmean{}}{\satsd{}}
  = \annsatzscore{}
\end{align*}

\begin{onebox}{Skor-Z}
  Skor-Z suatu pengamatan menyatakan berapa simpangan baku letak
  pengamatan tersebut di atas atau di bawah rata-rata.
  Kita menghitung skor-Z untuk pengamatan $x$ yang mengikuti distribusi
  dengan rata-rata $\mu$ dan simpangan baku $\sigma$ menggunakan
  \begin{align*}
  Z = \frac{x - \mu}{\sigma}
  \end{align*}
\end{onebox}

\begin{exercisewrap}
\begin{nexercise}
Gunakan skor ACT Tom, \tomsatscore{}, bersama rata-rata dan
simpangan baku ACT untuk mencari skor-Z-nya.\footnotemark{}
\end{nexercise}
\end{exercisewrap}
\footnotetext{$Z_{Tom}
  = \frac{x_{\text{Tom}} - \mu_{\text{ACT}}}
      {\sigma_{\text{ACT}}}
  = \frac{\tomsatscore{} - \actmean{}}{\actsd{}}
  = \tomsatzscore{}$}

Pengamatan di atas rata-rata selalu memiliki skor-Z positif,
sedangkan pengamatan di bawah rata-rata selalu memiliki skor-Z negatif.
Jika suatu pengamatan sama dengan rata-rata, misalnya skor SAT
sebesar \satmean{}, skor-Z-nya adalah $0$.

\begin{exercisewrap}
\begin{nexercise}
Misalkan $X$ merepresentasikan variabel acak dari
$N(\mu=3, \sigma=2)$ dan kita mengamati $x=5.19$. \\
%\begin{enumerate}[(a)]
%\setlength{\itemsep}{0mm}
%\item
(a)
    Carilah skor-Z untuk $x$. \\
%\item
(b)
    Gunakan skor-Z untuk menentukan berapa simpangan baku posisi $x$
    di atas atau di bawah rata-rata.\footnotemark{}
%\end{enumerate}
\end{nexercise}
\end{exercisewrap}
\footnotetext{(a) Skor-Z-nya diberikan oleh
    $Z
      = \frac{x-\mu}{\sigma}
      = \frac{5.19 - 3}{2}
      = 2.19/2
      = 1.095$.
    (b)~Pengamatan $x$ berada 1.095 simpangan baku
    \emph{di atas} rata-rata.
    Kita tahu pengamatan itu berada di atas rata-rata karena $Z$ positif.}

\begin{exercisewrap}
\begin{nexercise} \label{headLZScore}
Panjang kepala posum ekor sikat mengikuti distribusi normal
dengan rata-rata 92.6 mm dan simpangan baku 3.6 mm.
Hitung skor-Z untuk posum dengan panjang kepala 95.4 mm
dan 85.8~mm.\footnotemark{}
\end{nexercise}
\end{exercisewrap}
\footnotetext{Untuk $x_1=95.4$ mm:
    $Z_1
      = \frac{x_1 - \mu}{\sigma}
      = \frac{95.4 - 92.6}{3.6}
      = 0.78$.
    Untuk $x_2=85.8$ mm:
    $Z_2 = \frac{85.8 - 92.6}{3.6} = -1.89$.}

Kita dapat menggunakan skor-Z untuk mengenali secara kasar
pengamatan mana yang lebih tak lazim daripada yang lain.
Pengamatan $x_1$ disebut lebih tak lazim daripada pengamatan lain $x_2$
jika nilai mutlak skor-Z-nya lebih besar daripada nilai mutlak
skor-Z pengamatan lainnya: $|Z_1| > |Z_2|$.
Teknik ini sangat informatif ketika distribusinya simetris.

%\D{\newpage}

\begin{exercisewrap}
\begin{nexercise}
Manakah di antara pengamatan pada Latihan Terpandu~\ref{headLZScore}
yang lebih tak lazim?\footnotemark{}
\end{nexercise}
\end{exercisewrap}
\footnotetext{Karena \emph{nilai mutlak} skor-Z pengamatan kedua
  lebih besar daripada pengamatan pertama, pengamatan kedua
  memiliki panjang kepala yang lebih tak lazim.}

\subsection{Mencari luas ekor}

Dalam statistika, kemampuan menentukan luas ekor suatu distribusi
sangatlah berguna.
Sebagai contoh, berapa proporsi peserta yang memperoleh skor SAT
di bawah skor Ann sebesar 1300?
Nilai ini sama dengan \termsub{persentil}{percentile} Ann, yaitu
persentase peserta yang skornya lebih rendah daripada skor Ann.
Kita dapat memvisualisasikan luas ekor tersebut dengan kurva dan
arsiran seperti pada Gambar~\ref{satBelow1300}.

\begin{figure}[h]
  \centering
  \Figure[Ditampilkan sebuah distribusi normal dengan rata-rata 1100 dan simpangan baku 200. Distribusi diarsir di sebelah kiri nilai 1300; dengan kata lain, daerah yang dibatasi oleh sumbu horizontal, kurva berbentuk lonceng sampai nilai horizontal 1300, dan garis vertikal pada 1300 diberi arsiran.]{0.45}{satBelow1300}
  \caption{Luas di sebelah kiri $Z$ merepresentasikan proporsi peserta
      yang memperoleh skor lebih rendah daripada Ann.}
  \label{satBelow1300}
\end{figure}

Ada banyak teknik untuk melakukan perhitungan ini; kita akan membahas
tiga pilihan.
\begin{enumerate}
\item
    Pendekatan yang paling umum dalam praktik adalah menggunakan
    perangkat lunak statistika.
    Sebagai contoh, dalam program \R{}, kita dapat mencari luas
    pada Gambar~\ref{satBelow1300} dengan perintah berikut.
    Perintah ini menerima skor-Z dan mengembalikan luas ekor bawah: \\
    {\color{white}.....}%
        \texttt{> pnorm(1)} \\
    {\color{white}.....}%
        \texttt{[1] 0.8413447} \\
    Menurut perhitungan ini, daerah berarsir di bawah 1300
    merepresentasikan proporsi 0.841 (84.1\%) peserta SAT
    yang memiliki skor-Z di bawah $Z = 1$.
    Secara lebih umum, kita juga dapat menetapkan nilai batas secara
    eksplisit jika rata-rata dan simpangan bakunya turut diberikan: \\
    {\color{white}.....}%
        \texttt{> pnorm(1300, mean = 1100, sd = 200)} \\
    {\color{white}.....}%
        \texttt{[1] 0.8413447} %\\
    %\Add{More examples for using \R{} are provided
    %  at the end of the section.}


    Ada banyak pilihan perangkat lunak lain, seperti Python atau SAS;
    bahkan program lembar kerja seperti Excel dan Google Sheets
    mendukung perhitungan ini.
\item
    Strategi yang umum digunakan di kelas adalah memakai kalkulator
    grafik, misalnya kalkulator TI atau Casio.
    Kalkulator-kalkulator ini memerlukan serangkaian penekanan tombol
    yang tidak mudah dijelaskan secara ringkas.
    Petunjuk penggunaan kalkulator tersebut untuk mencari luas ekor
    distribusi normal tersedia di pustaka video OpenIntro:
    \begin{center}
    \oiRedirect{textbook-openintro_videos}
        {www.openintro.org/videos}
    \end{center}
\item
    Pilihan terakhir untuk mencari luas ekor adalah menggunakan
    \termsub{tabel peluang}{probability table}; tabel semacam ini
    kadang digunakan di kelas, tetapi jarang digunakan dalam praktik.
    Lampiran~\ref{normalProbabilityTable} memuat tabel tersebut
    beserta panduan penggunaannya.
\end{enumerate}
Dalam bagian ini, kita akan menyelesaikan soal distribusi normal
dengan selalu mencari skor-Z terlebih dahulu.
Alasannya, mulai Bab~\ref{ch_foundations_for_inf} kita akan menjumpai
konsep yang sangat serupa, yaitu
\indexthis{statistik uji}{test statistic}; dalam banyak keadaan,
statistik uji setara dengan skor-Z.

%No matter the approach you choose,
%try the Guided Practice exercises in this section
%using your preferred method.

\subsection{Contoh peluang normal}
\label{normal_probability_examples}

\noindent%
Skor total SAT dapat dihampiri dengan baik oleh model normal
$N(\mu = \satmean{}, \sigma = \satsd{})$.

\newcommand{\shannonsat}{1190}
\newcommand{\shannonsatz}{0.45}
\begin{examplewrap}
\begin{nexample}{Shannon dipilih secara acak dari para peserta SAT,
    dan tidak ada informasi tentang kemampuannya dalam SAT.
    Berapa peluang Shannon memperoleh skor sekurang-kurangnya
    \shannonsat{} pada SAT?}
  \label{satAbove1190Exam}%
  Pertama, selalu gambar dan beri keterangan pada kurva normal.
  (Gambar tidak harus tepat agar tetap berguna.)
  Kita tertarik pada peluang bahwa skornya melebihi \shannonsat{},
  sehingga kita mengarsir ekor atas berikut:
  \begin{center}
  \Figure[Ditampilkan sebuah distribusi normal dengan rata-rata 1100 dan simpangan baku 200. Daerah di bawah distribusi untuk nilai horizontal yang lebih besar dari 1190 diberi arsiran.]{0.4}{satAbove1190}
  \end{center}
  Gambar tersebut menampilkan rata-rata dan nilai-nilai yang terletak
  2~simpangan baku di atas dan di bawah rata-rata.
  Cara termudah untuk mencari luas berarsir di bawah kurva menggunakan
  skor-Z dari nilai batas.
  Dengan $\mu = \satmean{}$, $\sigma = \satsd{}$,
  dan nilai batas $x = \shannonsat{}$, skor-Z dihitung sebagai
  \begin{align*}
  Z = \frac{x - \mu}{\sigma}
    = \frac{\shannonsat{} - \satmean{}}{\satsd{}}
    = \frac{90}{\satsd{}}
    = \shannonsatz{}
  \end{align*}
  Dengan perangkat lunak statistika (atau metode lain yang dipilih),
  kita memperoleh luas di sebelah kiri $Z = \shannonsatz{}$ sebesar 0.6736.
  %This is Shannon's \term{percentile},
  %which is the fraction of folks who scored below her score
  %of \shannonsat{}.
  Untuk mencari luas \emph{di atas} $Z = \shannonsatz{}$,
  kita menghitung satu dikurangi luas ekor bawah:
  \begin{center}
  \Figure[Ditampilkan sebuah distribusi normal yang seluruh daerahnya diarsir, lalu tanda “dikurangi”, kemudian distribusi normal yang sebagian besar daerahnya diarsir sampai sedikit di atas rata-rata, lalu tanda sama dengan, dan akhirnya distribusi normal dengan daerah ekor atas diarsir. Di atas gambar-gambar tersebut tertulis “1.0000 dikurangi 0.6736 sama dengan 0.3264”. Visualisasi ini menunjukkan bahwa luas ekor atas distribusi normal dapat dipandang sebagai seluruh luas di bawah distribusi (yang bernilai 1) dikurangi luas di sebelah kiri sehingga tersisa luas di sebelah kanan.]{0.4}{subtractingArea}
  \end{center}
  Peluang Shannon memperoleh skor SAT sekurang-kurangnya 1190
  adalah 0.3264.
\end{nexample}
\end{examplewrap}

\begin{onebox}{Selalu gambar terlebih dahulu,
    lalu cari skor-Z}
  Dalam setiap situasi peluang normal,
  \emph{selalu, selalu, selalu} gambar dan beri keterangan pada
  kurva normal, kemudian arsir daerah yang ditanyakan.
  Gambar tersebut akan memberikan perkiraan peluang.
  Setelah menggambar situasinya, tentukan skor-Z untuk nilai
  yang ditanyakan.
\end{onebox}

\begin{exercisewrap}
\begin{nexercise}
Jika peluang Shannon memperoleh skor sekurang-kurangnya \shannonsat{}
adalah 0.3264, berapa peluang ia memperoleh skor kurang dari
\shannonsat{}?
Gambarlah kurva normal untuk latihan ini dan arsir daerah bawah,
bukan daerah atas.\footnotemark{}
\end{nexercise}
\end{exercisewrap}
\footnotetext{Kita telah memperoleh peluang ini pada
  Contoh~\ref{satAbove1190Exam}: 0.6736. \\
  \Figures[Sebuah distribusi normal dengan rata-rata 1100 dan simpangan baku 200 diarsir dari kiri sampai sebuah garis vertikal yang terletak sedikit di atas rata-rata distribusi.]{0.35}{subtractingArea}{subtracted}}

\D{\newpage}

\newcommand{\edwardsat}{1030}
\newcommand{\edwardsatz}{-0.35}
\newcommand{\edwardsatlower}{0.3632}
\begin{examplewrap}
\begin{nexample}{Edward memperoleh skor SAT sebesar \edwardsat{}.
    Berapa persentilnya?}
  \label{edwardSatBelow\edwardsat{}}%
  Pertama, kita perlu menggambar situasinya.
  \indexthis{Persentil}{percentile} Edward adalah proporsi peserta
  yang tidak mencapai skor setinggi \edwardsat{}.
  Skor-skor ini berada di sebelah kiri \edwardsat{}.
\begin{center}
\Figure[Distribusi normal dengan rata-rata 1100 dan simpangan baku 200 diarsir dari kiri sampai sebuah garis vertikal yang terletak sedikit di bawah rata-rata distribusi, pada skor 1030.]{0.3}{satBelow1030}
\end{center}
Dengan menentukan rata-rata $\mu=\satmean{}$, simpangan baku
$\sigma=\satsd{}$, dan nilai batas untuk luas ekor $x=\edwardsat{}$,
kita dapat menghitung skor-Z dengan mudah:
\begin{align*}
Z
  = \frac{x - \mu}{\sigma}
  = \frac{\edwardsat{} - \satmean{}}{\satsd{}}
  = \edwardsatz{}
\end{align*}
Dengan perangkat lunak statistika, kita memperoleh luas ekor 0.3632.
Edward berada pada persentil ke-$36$.
\end{nexample}
\end{examplewrap}

\begin{exercisewrap}
\begin{nexercise}
Gunakan hasil Contoh~\ref{edwardSatBelow\edwardsat{}}
untuk menghitung proporsi peserta SAT yang memperoleh skor
lebih tinggi daripada Edward.
Gambarlah pula situasi yang baru.\footnotemark{}
\end{nexercise}
\end{exercisewrap}
\footnotetext{Jika Edward memperoleh skor lebih tinggi daripada 36\%
  peserta SAT, sekitar 64\% peserta tentu memperoleh skor lebih tinggi
  daripada Edward. \\
  \Figures[Ditampilkan dua distribusi normal dengan rata-rata 1100 dan simpangan baku 200. Distribusi pertama diarsir di sebelah kiri skor 1030, sedangkan distribusi kedua diarsir di sebelah kanan skor 1030.]{0.25}{satBelow1030}{satAbove1030}}

\begin{onebox}{Mencari luas di sebelah kanan}
  Ketika diberi skor-Z, banyak program perangkat lunak mengembalikan
  luas di sebelah kiri.
  Jika Anda memerlukan luas di sebelah kanan, cari dahulu luas
  di sebelah kiri lalu kurangi nilai tersebut dari satu.
\end{onebox}

\newcommand{\stuartsat}{1500}
\newcommand{\stuarsatz}{2}
\begin{exercisewrap}
\begin{nexercise}
Stuart memperoleh skor SAT sebesar \stuartsat{}.
Gambarlah situasi untuk setiap bagian. \\
(a)~Berapa persentilnya? \\
(b)~Berapa persen peserta SAT yang memperoleh skor lebih tinggi
  daripada Stuart?\footnotemark{}
\end{nexercise}
\end{exercisewrap}
\footnotetext{Gambarnya kami serahkan kepada Anda.
  (a) $Z = \frac{\stuartsat{} - \satmean{}}{\satsd{}}
         = \stuarsatz{}
         \to 0.9772$.
  (b) $1 - 0.9772 = 0.0228$.}

Berdasarkan sampel 100 laki-laki, tinggi badan laki-laki dewasa
di AS hampir berdistribusi normal dengan rata-rata 70.0''
dan simpangan baku 3.3''.

\begin{exercisewrap}
\begin{nexercise}
Mike memiliki tinggi 5'7'' dan Jose 6'4'', dan keduanya tinggal di AS. \\
(a) Berapa persentil tinggi badan Mike? \\
(b) Berapa persentil tinggi badan Jose? \\
Gambarlah pula satu situasi untuk setiap bagian.\footnotemark{}
\end{nexercise}
\end{exercisewrap}
\footnotetext{Ubah dahulu tinggi badan menjadi inci:
  67 dan 76 inci.
  Gambar-gambarnya ditampilkan di bawah. \\
  (a) $Z_{\text{Mike}} = \frac{67 - 70}{3.3} = -0.91\ \to\ 0.1814$.
  (b) $Z_{\text{Jose}} = \frac{76 - 70}{3.3} = 1.82\ \to\ 0.9656$.
  \\
  \Figure[Ditampilkan dua grafik. Grafik pertama diberi label “Mike” dan menampilkan distribusi normal dengan rata-rata 70 serta ekor kiri di bawah 67 diarsir. Grafik kedua diberi label “Jose” dan menampilkan distribusi normal dengan rata-rata 70 serta sebagian besar distribusi sampai nilai 76 diarsir.]{0.45}{mikeAndJosePercentiles}}

\D{\newpage}

Beberapa soal terakhir berfokus pada pencarian persentil (ekor bawah)
atau ekor atas untuk suatu pengamatan tertentu.
Bagaimana jika Anda ingin mengetahui pengamatan yang bersesuaian
dengan persentil tertentu?

\begin{examplewrap}
\begin{nexample}{Tinggi Erik berada pada persentil ke-$40$.
    Berapakah tinggi Erik?}\label{normalExam40Perc}
  Seperti biasa, gambar dahulu situasinya.\vspace{-4mm}
  \begin{center}
  \Figure[Ditampilkan sebuah distribusi normal untuk tinggi badan laki-laki dewasa dengan rata-rata 70 inci dan simpangan baku 3.3 inci. Daerah dari kiri sampai suatu nilai sedikit di bawah rata-rata diarsir dan diberi label 40\% (0.40).]{0.3}{height40Perc}\vspace{-1mm}
  \end{center}
  Dalam kasus ini, peluang ekor bawah diketahui (0.40) dan dapat
  diarsir pada diagram.
  Kita ingin mencari pengamatan yang bersesuaian dengan nilai tersebut.
  Sebagai langkah awal, kita menentukan skor-Z yang bersesuaian
  dengan persentil ke-$40$.
  Dengan perangkat lunak, kita memperoleh skor-Z yang bersesuaian,
  yaitu sekitar -0.25.

  Setelah mengetahui $Z_{Erik} = -0.25$ serta parameter populasi
  $\mu = 70$ dan $\sigma = 3.3$ inci, kita dapat menyusun rumus skor-Z
  untuk menentukan tinggi Erik yang belum diketahui, yang diberi label
  $x_{_{\text{Erik}}}$:
  \begin{align*}
  -0.25
    = Z_{_{\text{Erik}}}
    = \frac{x_{_{\text{Erik}}} - \mu}{\sigma}
    = \frac{x_{_{\text{Erik}}} - 70}{3.3}
  \end{align*}
  Dengan menyelesaikan persamaan untuk $x_{_{\text{Erik}}}$,
  diperoleh tinggi 69.18 inci.
  Jadi, tinggi Erik sekitar 5'9''.
\end{nexample}
\end{examplewrap}

\begin{examplewrap}
\begin{nexample}{Berapa tinggi laki-laki dewasa pada
    persentil ke-$82$?}
  Sekali lagi, kita menggambar situasinya terlebih dahulu.\vspace{-3mm}
  \begin{center}
  \Figure[Distribusi normal dengan rata-rata 70 dan simpangan baku 3.3 diarsir dari kiri sampai sebuah garis vertikal yang terletak sedikit di atas rata-rata distribusi. Luas berarsir di sebelah kiri garis vertikal diberi label “82\% (0.82)” dan ekor atas yang tidak diarsir diberi label “18\% (0.18)”.]{0.28}{height82Perc}\vspace{-1mm}
  \end{center}
  Berikutnya, kita ingin mencari skor-Z pada persentil ke-$82$.
  Nilainya positif dan dapat diperoleh dengan perangkat lunak sebagai
  $Z = 0.92$.
  Terakhir, tinggi $x$ dicari menggunakan rumus skor-Z dengan rata-rata
  $\mu$, simpangan baku $\sigma$, dan skor-Z $Z = 0.92$ yang diketahui:
  \begin{align*}
  0.92 = Z = \frac{x-\mu}{\sigma} = \frac{x - 70}{3.3}
  \end{align*}
  Hasilnya adalah 73.04 inci, atau sekitar 6'1'', sebagai tinggi
  pada persentil ke-$82$.
\end{nexample}
\end{examplewrap}

\begin{exercisewrap}
\begin{nexercise}
Skor SAT mengikuti $N(\satmean{}, \satsd{})$.\footnotemark{} \\
(a) Berapakah skor SAT pada persentil ke-$95$? \\
(b) Berapakah skor SAT pada persentil ke-$97.5$?
\end{nexercise}
\end{exercisewrap}
\footnotetext{Jawaban singkat:
  (a) $Z_{95} = 1.6449 \to 1429$ skor SAT.
  (b) $Z_{97.5} = 1.96 \to 1492$ skor SAT.}

\D{\newpage}

\begin{exercisewrap}
\begin{nexercise}\label{more74Less69}
Tinggi badan laki-laki dewasa mengikuti $N(70.0$''$, 3.3$''$)$.\footnotemark{} \\
(a)~Berapa peluang bahwa laki-laki dewasa yang dipilih secara acak
    memiliki tinggi sekurang-kurangnya 6'2'' (74 inci)? \\
(b)~Berapa peluang bahwa seorang laki-laki dewasa memiliki tinggi
    kurang dari 5'9'' (69 inci)?
\end{nexercise}
\end{exercisewrap}
\footnotetext{Jawaban singkat:
  (a) $Z = 1.21 \to 0.8869$, kemudian kurangi nilai ini
      dari 1 untuk memperoleh 0.1131.
  (b) $Z = -0.30 \to 0.3821$.}

\begin{examplewrap}
\begin{nexample}{Berapa peluang bahwa laki-laki dewasa yang dipilih
    secara acak memiliki tinggi antara 5'9'' dan 6'2''?}
  Tinggi tersebut bersesuaian dengan 69 inci dan 74 inci.
  Pertama, gambar situasinya.
  Daerah yang ditanyakan bukan lagi ekor atas atau ekor bawah.\vspace{-2mm}
  \begin{center}
  \Figure[Ditampilkan distribusi normal dengan rata-rata 70 dan simpangan baku 3.3. Daerah dari nilai yang sedikit di bawah rata-rata (69) sampai nilai yang lebih jauh pada ekor kanan (74) diberi arsiran.]{0.35}{between59And62}\vspace{-2mm}
  \end{center}
  Luas total di bawah kurva adalah~1.
  Jika kita mencari luas kedua ekor yang tidak diarsir
  (dari Latihan Terpandu~\ref{more74Less69}, luasnya masing-masing
  $0.3821$ dan $0.1131$), kita dapat mencari luas bagian tengah:
  \vspace{-2mm}
  \begin{center}
  \Figure[Ditampilkan perhitungan berbentuk grafik: seluruh distribusi (1.0000) dikurangi ekor bawah (0.3821) dan ekor atas kecil (0.1131), sehingga tersisa distribusi normal dengan hanya satu ruas berarsir, dari sedikit di bawah rata-rata sampai agak di atas rata-rata. Luas berarsir terakhir diberi label 0.5048.]{0.55}{subtracting2Areas}\vspace{-2mm}
  \end{center}
  Jadi, peluang tinggi badan berada antara 5'9'' dan 6'2''
  adalah 0.5048.
\end{nexample}
\end{examplewrap}

\begin{exercisewrap}
\begin{nexercise}
Skor SAT mengikuti $N(\satmean{}, \satsd{})$.
Berapa persen peserta SAT yang memperoleh skor antara \satmean{}
dan 1400?\footnotemark
\end{nexercise}
\end{exercisewrap}
\footnotetext{Berikut solusi ringkasnya.
  (Pastikan Anda menggambar situasinya!)
  Pertama, cari persentase peserta yang memperoleh skor di bawah
  \satmean{} dan persentase yang memperoleh skor di atas 1400:
  $Z_{\satmean{}} = 0.00 \to 0.5000$ (luas di bawah),
  $Z_{1400} = 1.5 \to 0.0668$ (luas di atas).
  Jawaban akhir: $1.0000 - 0.5000 - 0.0668 = 0.4332$.}

\begin{exercisewrap}
\begin{nexercise}
Tinggi badan laki-laki dewasa mengikuti $N(70.0$''$, 3.3$''$)$.
Berapa persen laki-laki dewasa yang tingginya antara 5'5''
dan 5'7''?\footnotemark{}
\end{nexercise}
\end{exercisewrap}
\footnotetext{5'5'' adalah 65 inci ($Z = -1.52$).
  5'7'' adalah 67 inci ($Z = -0.91$).
  Solusi numerik: $1.000 - 0.0643 - 0.8186 = 0.1171$,
  yaitu 11.71\%.}

\D{\newpage}

\subsection{Kaidah 68-95-99.7}

Di sini kita memperkenalkan kaidah praktis untuk peluang bahwa
suatu nilai berada dalam jarak 1, 2, dan 3 simpangan baku dari
rata-rata pada distribusi normal.
Kaidah ini berguna dalam beragam situasi praktis, terutama ketika
kita ingin membuat perkiraan cepat tanpa kalkulator atau tabel-Z.

\begin{figure}[hht]
\centering
\Figure[Ditampilkan sebuah distribusi normal. Daerah tengah, dari satu simpangan baku di bawah rata-rata sampai satu simpangan baku di atas rata-rata, diarsir biru dan diberi label 68\%. Daerah yang lebih jauh, dari dua simpangan baku di bawah rata-rata sampai dua simpangan baku di atas rata-rata, diarsir hijau selain bagian yang telah diarsir biru dan diberi label 95\%. Daerah yang lebih jauh lagi, dari tiga simpangan baku di bawah rata-rata sampai tiga simpangan baku di atas rata-rata, diarsir kuning selain bagian yang telah diarsir hijau atau biru dan diberi label 99.7\%. Persentase 68\%, 95\%, dan 99.7\% tersebut merepresentasikan bagian luas di bawah distribusi normal yang terletak dalam 1, 2, dan 3 simpangan baku dari rata-rata.]{0.63}{6895997}
\caption{Peluang berada dalam jarak 1, 2, dan 3 simpangan baku
  dari rata-rata pada distribusi normal.}
\label{6895997}
\end{figure}

\begin{exercisewrap}
\begin{nexercise}
Gunakan perangkat lunak, kalkulator, atau tabel peluang untuk
mengonfirmasi bahwa sekitar 68\%, 95\%, dan 99.7\% pengamatan
masing-masing berada dalam jarak 1, 2, dan 3 simpangan baku
dari rata-rata pada distribusi normal.
Sebagai contoh, carilah dahulu luas antara $Z=-1$ dan $Z=1$,
yang seharusnya sekitar 0.68.
Demikian pula, luas antara $Z=-2$ dan $Z=2$ seharusnya sekitar 0.95.
\footnotemark{}
\end{nexercise}
\end{exercisewrap}
\footnotetext{Gambarlah situasinya terlebih dahulu.
  Dengan perangkat lunak, kita memperoleh 0.6827 dalam 1~simpangan baku,
  0.9545 dalam 2~simpangan baku,
  dan 0.9973 dalam 3~simpangan baku.}

Variabel acak normal mungkin saja berada 4,~5, atau bahkan lebih banyak
simpangan baku dari rata-rata.
Namun, kejadian semacam ini sangat jarang jika datanya hampir normal.
Peluang berada lebih dari 4 simpangan baku dari rata-rata adalah
sekitar 1 dari 15,000.
Untuk 5 dan 6 simpangan baku, peluangnya masing-masing sekitar
1 dari 2 juta dan 1 dari 500 juta.

\begin{exercisewrap}
\begin{nexercise}
Distribusi skor SAT dapat dihampiri dengan baik oleh model normal
dengan rata-rata $\mu = \satmean{}$ dan simpangan baku
$\sigma = \satsd{}$.\footnotemark{} \\
(a) Kira-kira berapa persen peserta ujian yang memperoleh skor
    antara 700 dan 1500? \\
(b) Berapa persen yang memperoleh skor antara \satmean{} dan 1500?
\end{nexercise}
\end{exercisewrap}
\footnotetext{(a) 700 dan 1500 berada dua simpangan baku
  di bawah dan di atas rata-rata, sehingga sekitar 95\% peserta ujian
  memperoleh skor antara 700 dan 1500.
  (b)~Kita telah menemukan bahwa rentang 700 sampai 1500 mencakup
  sekitar 95\% peserta ujian.
  Para peserta ini terbagi sama banyak oleh pusat distribusi,
  yaitu \satmean{}, sehingga
  $\frac{95\%}{2} = 47.5\%$ dari seluruh peserta ujian
  memperoleh skor antara \satmean{} dan 1500.}


{\input{ch_distributions/TeX/normal_distribution.tex}}
